% Chapters/5-QED-for-leptons.tex

\chapter{QED for Leptons}
\label{ch:qed_leptons}

\section{QED for elementary spin-$\tfrac12$ particles}
Quantum Electrodynamics (QED) is the relativistic quantum field theory of a $U(1)$ gauge field
(the photon) coupled to electrically charged matter. For a single elementary Dirac particle
of charge $q$ (and its antiparticle), the QED Lagrangian reads
\begin{equation}
\mathcal{L}_{\rm QED}
= -\frac14 F_{\mu\nu}F^{\mu\nu} + \bar\psi\left(i\gamma^\mu D_\mu - m\right)\psi,
\qquad
F_{\mu\nu}=\partial_\mu A_\nu-\partial_\nu A_\mu,
\qquad
D_\mu=\partial_\mu+i q A_\mu.
\end{equation}
At a given energy scale, the notion of ``elementary'' depends on what degrees of freedom can be
resolved. In this chapter we focus on leptons as elementary particles.

If we have $n$ elementary charged leptons $\psi_j$ with charges $q_j$ and masses $m_j$, then
\begin{equation}
\mathcal{L}_{\rm QED}
= -\frac14 F_{\mu\nu}F^{\mu\nu}
+ \sum_{j=1}^{n}\bar\psi_j\left(i\gamma^\mu D^{(j)}_\mu - m_j\right)\psi_j,
\qquad
D^{(j)}_\mu=\partial_\mu+i q_j A_\mu.
\end{equation}

\subsection{Global symmetries and flavor conservation}
The Lagrangian is invariant under \emph{independent global} phase rotations for each flavor,
\begin{equation}
\psi_j \to e^{i\theta_j}\psi_j,
\qquad \theta_j=\text{constant},
\end{equation}
so there are $n$ conserved charges. Physically: \textbf{each lepton flavor number is conserved in QED}
(e.g.\ electron number, muon number, etc.).

\subsection{Interaction Lagrangian and electromagnetic current}
Expanding the covariant derivatives shows that the interaction is
\begin{equation}
\mathcal{L}_I = \sum_{j=1}^n \left(-q_j \bar\psi_j\gamma^\mu A_\mu \psi_j\right)
\equiv -\sum_{j=1}^n j^\mu_j A_\mu \equiv - j^\mu A_\mu,
\qquad
j^\mu_j = q_j\,\bar\psi_j\gamma^\mu\psi_j.
\end{equation}
Most elementary QED processes start at second order in $\mathcal{L}_I$ (one virtual photon exchange).

\section{Lepton--lepton scattering: $e^-\mu^- \to e^-\mu^-$}
Two historically important processes (for the development of QCD) are
\begin{itemize}
  \item elastic scattering: $e^-\mu^- \to e^-\mu^-$,
  \item pair creation: $e^+e^- \to \mu^+\mu^-$,
\end{itemize}
which are related by \emph{crossing}. We begin with $e^-\mu^- \to e^-\mu^-$.

\subsection{Amplitude at second order in $\mathcal{L}_I$}
At leading order (one-photon exchange),
\begin{equation}
i\mathcal{M}
=
\frac{i^2}{2!}\int d^4x\;
\langle f|T\{\mathcal{L}_I(0)\mathcal{L}_I(x)\}|i\rangle.
\end{equation}
Writing $\mathcal{L}_I=-j^\mu A_\mu$ and separating electron and muon currents,
\begin{equation}
i\mathcal{M}
=
\frac{i^2}{2}\int d^4x\;
\langle f|T\{j_e^\mu(0)A_\mu(0)\; j_\mu^\nu(x)A_\nu(x)\}|i\rangle.
\end{equation}
Using Wick contractions for the photon field gives the photon propagator,
\begin{equation}
\langle 0|T\{A_\mu(0)A_\nu(x)\}|0\rangle
=
\int\frac{d^4k}{(2\pi)^4}\;
\frac{-i g_{\mu\nu}}{k^2+i\eta}\;e^{-ik\cdot(0-x)},
\end{equation}
so the amplitude factorizes into a propagator times current matrix elements.

\subsection{External states and tree-level result}
Let the initial and final lepton states be
\begin{equation}
|i\rangle = |e^-(p_A,\lambda_A)\rangle\;|\mu^-(p_B,\lambda_B)\rangle,
\qquad
|f\rangle = |e^-(p_1,\lambda_1)\rangle\;|\mu^-(p_2,\lambda_2)\rangle.
\end{equation}
The electron and muon currents between one-particle states are
\begin{equation}
\langle e^-(p_1)|\,j^\mu_e(0)\,|e^-(p_A)\rangle
= q_e\,\bar u(p_1)\gamma^\mu u(p_A),
\qquad
\langle \mu^-(p_2)|\,j^\nu_\mu(x)\,|\mu^-(p_B)\rangle
= q_\mu\,\bar u(p_2)\gamma^\nu u(p_B)\,e^{-ix\cdot(p_B-p_2)}.
\end{equation}
Performing the $x$-integral fixes the photon momentum $q\equiv p_B-p_2=p_1-p_A$ and yields
\begin{equation}
i\mathcal{M}
=
\frac{-i g_{\mu\nu}}{q^2}\;
\bar u(p_1)\,(-iq_e\gamma^\mu)\,u(p_A)\;
\bar u(p_2)\,(-iq_\mu\gamma^\nu)\,u(p_B).
\end{equation}
For $e^-$ and $\mu^-$, $q_e=q_\mu=-e$, hence (up to an irrelevant overall phase)
\begin{equation}
\mathcal{M}
=
e^2\;\frac{\bar u(p_1)\gamma^\mu u(p_A)\;\bar u(p_2)\gamma_\mu u(p_B)}{q^2}.
\end{equation}
In Mandelstam variables, $q^2=t$ (for elastic scattering), so
\begin{equation}
\mathcal{M}
=
e^2\;\frac{\bar u(p_1)\gamma^\mu u(p_A)\;\bar u(p_2)\gamma_\mu u(p_B)}{t}.
\end{equation}

\subsection{Spin averages, traces, and $|\mathcal{M}|^2$}
If the initial beams are unpolarized and final spins are not measured, we average/sum:
\begin{equation}
|\mathcal{M}|^2 \;\to\;
\overline{|\mathcal{M}|^2}
=
\frac12\sum_{\lambda_A=\pm}\frac12\sum_{\lambda_B=\pm}
\sum_{\lambda_1=\pm}\sum_{\lambda_2=\pm}
|\mathcal{M}|^2.
\end{equation}
Using completeness relations and trace technology,
\begin{align}
\sum_{\lambda_A,\lambda_1}\bar u(p_1)\gamma^\mu u(p_A)\;
\left(\bar u(p_1)\gamma^\nu u(p_A)\right)^\ast
&=
\mathrm{tr}\!\left[\gamma^\mu(\slashed{p}_A+m_e)\gamma^\nu(\slashed{p}_1+m_e)\right]
\equiv 2L^{\mu\nu}_e,\\
\sum_{\lambda_B,\lambda_2}\bar u(p_2)\gamma^\mu u(p_B)\;
\left(\bar u(p_2)\gamma^\nu u(p_B)\right)^\ast
&=
\mathrm{tr}\!\left[\gamma^\mu(\slashed{p}_B+m_\mu)\gamma^\nu(\slashed{p}_2+m_\mu)\right]
\equiv 2L^{\mu\nu}_\mu,
\end{align}
with
\begin{equation}
L^{\mu\nu}_e
=
2\left(p_A^\mu p_1^\nu + p_A^\nu p_1^\mu - g^{\mu\nu}\,p_A\!\cdot p_1\right)
+2m_e^2 g^{\mu\nu},
\qquad
L^{\mu\nu}_\mu
=
2\left(p_B^\mu p_2^\nu + p_B^\nu p_2^\mu - g^{\mu\nu}\,p_B\!\cdot p_2\right)
+2m_\mu^2 g^{\mu\nu}.
\end{equation}
Then
\begin{equation}
\overline{|\mathcal{M}|^2}
=
\frac{e^4}{t^2}\;L^{\mu\nu}_e L_{\mu\nu}^{\mu}.
\end{equation}
It is convenient to express the result in terms of Mandelstam variables
\begin{equation}
s\equiv (p_A+p_B)^2=(p_1+p_2)^2,\qquad
t\equiv (p_A-p_1)^2=(p_2-p_B)^2,\qquad
u\equiv (p_A-p_2)^2=(p_1-p_B)^2,
\end{equation}
satisfying $s+t+u=2m_e^2+2m_\mu^2$. One finds
\begin{equation}
L^{\mu\nu}_e L_{\mu\nu}^{\mu}
=
4\left[(s-m_e^2-m_\mu^2)^2+\frac{t^2}{2}+st\right],
\end{equation}
hence
\begin{equation}
\overline{|\mathcal{M}|^2}
=
\frac{4e^4}{t^2}\left[(s-m_e^2-m_\mu^2)^2+\frac{t^2}{2}+st\right]
\simeq
\frac{4e^4}{t^2}\left[(s-m_\mu^2)^2+\frac{t^2}{2}+st\right]
\quad (m_\mu\gg m_e).
\end{equation}

\subsection{Differential cross section in the CoM frame}
For a $2\to 2$ process in the center-of-mass frame,
\begin{equation}
\left(\frac{d\sigma}{d\Omega}\right)_{\rm CoM}
=
\frac{\overline{|\mathcal{M}|^2}}{64\pi^2 s}\;\frac{p_f}{p_i},
\end{equation}
with $p_i=|\vec p_A|=|\vec p_B|$ and $p_f=|\vec p_1|=|\vec p_2|$.
For $e^-\mu^-\to e^-\mu^-$ we have $p_f=p_i$.

In the \textbf{high-energy limit} ($p_i\gg m_\mu$),
\begin{equation}
s\simeq 4p_i^2,
\qquad
t\simeq -4p_i^2\sin^2\frac{\theta}{2},
\qquad
u\simeq -4p_i^2\cos^2\frac{\theta}{2},
\end{equation}
and the squared amplitude becomes
\begin{equation}
\overline{|\mathcal{M}|^2}
\simeq
2e^4\;\frac{s^2+u^2}{t^2}
\simeq
2e^4\;\frac{1+\cos^4(\theta/2)}{\sin^4(\theta/2)}.
\end{equation}
This is independent of $p_i$ and is strongly \textbf{forward enhanced} (diverges as $\theta\to 0$),
a classic signature of long-range Coulomb exchange.

\subsection{Mott and Rutherford limits}
In the regime $m_e\ll p_i\ll m_\mu$ (electron scattering off a heavy target in CoM kinematics),
one recovers Mott's formula (essentially relativistic Rutherford with spin effects),
\begin{equation}
\left(\frac{d\sigma}{d\Omega}\right)_{\rm CoM}
\simeq
\frac{\alpha^2}{4p_i^2}\;\frac{\cos^2(\theta/2)}{\sin^4(\theta/2)},
\qquad
\alpha\equiv \frac{e^2}{4\pi}.
\end{equation}
In the non-relativistic limit $p_i\ll m_e$ (so $E_A\simeq m_e$), Rutherford scattering is recovered:
\begin{equation}
\left(\frac{d\sigma}{d\Omega}\right)_{\rm CoM}
\simeq
\frac{\alpha^2 m_e^2}{4p_i^4}\;\frac{1}{\sin^4(\theta/2)}.
\end{equation}

\section{The LAB frame (fixed-target kinematics)}
The LAB frame is natural for fixed-target experiments. Take
\begin{equation}
p_A=(E_A,\vec p_A),\qquad p_B=(m_B,\vec 0),
\qquad
p_1=(E_1,\vec p_1),\qquad p_2=(E_2,\vec p_2),
\end{equation}
with $m_B=m_\mu$ for a muon target. Define the momentum transfer
\begin{equation}
q\equiv p_A-p_1,
\qquad
\nu\equiv \frac{q\cdot p_B}{m_B}=E_A-E_1=q^0.
\end{equation}
After integrating over the final target momentum, one finds a convenient expression where a delta function
enforces on-shell kinematics. In the limit $E_A\gg m_e$, the delta function fixes $E_1$ to
\begin{equation}
E_1=\frac{E_A}{1+\frac{2E_A}{m_B}\sin^2(\theta/2)}.
\end{equation}
The differential cross section can be written as
\begin{equation}
\left(\frac{d\sigma}{d\Omega}\right)_{\rm LAB}
=
\frac{\overline{|\mathcal{M}|^2}}{64\pi^2 m_B^2}\;
\frac{1}{\left(1+\frac{2E_A}{m_B}\sin^2(\theta/2)\right)^2}
=
\frac{\overline{|\mathcal{M}|^2}\,E_1^2}{64\pi^2 m_B^2 E_A^2},
\end{equation}
with $q^2=t\simeq -4E_AE_1\sin^2(\theta/2)$. In the limit $m_e\ll E_A\ll m_\mu$ one recovers the LAB version of Mott:
\begin{equation}
\left(\frac{d\sigma}{d\Omega}\right)_{\rm LAB}
\simeq
\frac{\alpha^2}{4E_A^2}\;\frac{\cos^2(\theta/2)}{\sin^4(\theta/2)}.
\end{equation}

\section{Pair creation: $e^+e^- \to \mu^+\mu^-$}
This process is related to $e^-\mu^-\to e^-\mu^-$ by \textbf{crossing}. For unpolarized beams and unmeasured final spins,
the squared amplitude is obtained by the substitution
\begin{equation}
s\to t,\qquad t\to s,\qquad u\to u,
\end{equation}
leading to
\begin{equation}
\overline{|\mathcal{M}|^2}
=
\frac{4e^4}{s^2}\left[(t-m_e^2-m_\mu^2)^2+\frac{s^2}{2}+st\right]
\simeq
\frac{4e^4}{s^2}\left[(t-m_\mu^2)^2+\frac{s^2}{2}+st\right].
\end{equation}

\subsection{High-energy limit in the CoM frame}
For $s\gg m_\mu^2$ and in the CoM frame,
\begin{equation}
s\simeq 4E_A^2,
\qquad
t\simeq -2E_A^2(1-\cos\theta),
\qquad
u\simeq -2E_A^2(1+\cos\theta),
\end{equation}
one finds the classic result
\begin{equation}
\overline{|\mathcal{M}|^2}_{\rm CoM}
\simeq
e^4\left(1+\cos^2\theta\right),
\qquad
\left(\frac{d\sigma}{d\Omega}\right)_{\rm CoM}
\simeq
\frac{\alpha^2}{4s}\left(1+\cos^2\theta\right).
\end{equation}
Integrating over angles gives
\begin{equation}
\sigma_{\rm CoM}
=
\frac{4\pi\alpha^2}{3s}.
\end{equation}
The angular distribution peaks at $\theta=0,\pi$.

\subsection{Near threshold: $\beta$-factor}
Close to threshold ($s\gtrsim 4m_\mu^2$), the outgoing muons are slow. Define
\begin{equation}
\beta \equiv \sqrt{1-\frac{4m_\mu^2}{s}},
\qquad
p_f = \frac{\sqrt{s}}{2}\beta.
\end{equation}
The result becomes (showing the leading mass correction structure)
\begin{equation}
\left(\frac{d\sigma}{d\Omega}\right)_{\rm CoM}
\simeq
\frac{\alpha^2\beta}{4s}\left[(1+\cos^2\theta)+\frac{4m_\mu^2}{s}\sin^2\theta\right],
\qquad
\sigma_{\rm CoM}
\simeq
\frac{4\pi\alpha^2\beta}{3s}\left(1+\frac{2m_\mu^2}{s}\right).
\end{equation}

\section{Scalar pair production: $e^+e^- \to \pi^+\pi^-$ at low energy}
At energies $E_A\ll 1~\text{GeV}$ one may treat pions approximately as point-like scalars.
Then one uses QED for the electron plus \textbf{scalar QED} for the pion field $\phi$.
The relevant interactions are
\begin{equation}
\mathcal{L}_I
=
-q_e\,\bar\psi\gamma^\mu A_\mu\psi
+ i q_\pi A^\mu\left(\partial_\mu\phi^\ast\,\phi-\phi^\ast\,\partial_\mu\phi\right)
+ q_\pi^2 A_\mu A^\mu\,\phi^\ast\phi.
\end{equation}
The last term is quadratic in $q_\pi$ and does not contribute at leading order for one-photon exchange.

Proceeding as before, one finds (for $E_A\gg m_e$) an angular distribution in the CoM frame of the form
\begin{equation}
\overline{|\mathcal{M}|^2}_{\rm CoM}
\simeq
e^4\,\beta^2\sin^2\theta,
\qquad
\beta \equiv \sqrt{1-\frac{4m_\pi^2}{s}}.
\end{equation}
Hence,
\begin{equation}
\left(\frac{d\sigma}{d\Omega}\right)_{\rm CoM}
\simeq
\frac{\alpha^2\beta^3}{8s}\,\sin^2\theta,
\qquad
\sigma_{\rm CoM}
\simeq
\frac{\pi\alpha^2\beta^3}{3s}.
\end{equation}
\textbf{Key physics:} the angular distribution peaks at $\theta=\pi/2$ for scalar production,
whereas $\mu^+\mu^-$ production peaks at $\theta=0,\pi$. Measuring angular distributions allows
one to infer the spin of the produced particles.

For $s\gg m_\pi^2$, one gets $\sigma_{\pi\pi}\simeq \pi\alpha^2/(3s)$, i.e.\ a factor of $1/4$
relative to $\sigma_{\mu\mu}=4\pi\alpha^2/(3s)$.

\section{Other elementary QED processes (overview)}
Besides lepton--lepton scattering and pair creation, several $2\to2$ reactions appear repeatedly:

\subsection{M{\o}ller scattering: $e^-e^- \to e^-e^-$}
\begin{itemize}
  \item Two tree-level diagrams contribute (exchange of identical fermions), with a \textbf{relative minus sign}.
  \item The final state contains two identical particles, so the cross section must be divided by $2!$.
  \item Related by crossing to Bhabha scattering and to $e^+e^+\to e^+e^+$.
\end{itemize}

\subsection{Bhabha scattering: $e^-e^+ \to e^-e^+$}
\begin{itemize}
  \item Two tree-level diagrams contribute (annihilation and scattering channels), with a \textbf{relative minus sign}.
  \item Related by crossing to M{\o}ller scattering.
\end{itemize}

\subsection{Compton scattering: $e^-\gamma \to e^-\gamma$}
\begin{itemize}
  \item Two diagrams contribute (the photon attaches before/after the interaction along the electron line), with the \textbf{same sign}.
  \item Photon polarization vectors enter explicitly in the amplitude.
  \item The internal line is a \textbf{Dirac propagator}.
\end{itemize}

\subsection{Annihilation into photons: $e^-e^+ \to \gamma\gamma$}
\begin{itemize}
  \item Two diagrams contribute, with the \textbf{same sign}.
  \item Photon polarizations enter explicitly.
  \item Related by crossing to Compton scattering and to $\gamma\gamma\to e^-e^+$.
  \item Important for positronium ($e^-e^+$ bound state) decay.
\end{itemize}

\section{Summary}
QED provides a clean laboratory where:
\begin{itemize}
  \item global flavor symmetries imply lepton flavor conservation (in pure QED),
  \item tree-level amplitudes follow from one-photon exchange and factorize into currents,
  \item spin sums reduce to gamma-matrix traces and lead naturally to expressions in Mandelstam variables,
  \item angular distributions and threshold factors ($\beta$) encode the spin and mass of produced particles.
\end{itemize}
These techniques will be essential when moving to non-abelian gauge theories and to processes
involving hadrons.